\chapter{Reproducibility}
\label{app:repro}

The code is in the project repository. The all-mixed dataset \cite{bnn2023allmixed} is expected under \path{data/}: the original splits in \path{data/all_mixed/} (used by the baseline) and the re-split in \path{data/all_mixed_v1/} (used by the teacher and students). The re-split validation set contains the topologies \path{topo-10-7}, \path{topo-50-0}, \path{topo-75-10}, \path{topo-100-20}, \path{topo-130-26} and \path{topo-170-34}.

\begin{table}[H]
\centering
\footnotesize
\caption{Main scripts, selected checkpoints and result files.}
\label{tab:repro}
\begin{tabularx}{\textwidth}{@{}p{22mm} Y >{\raggedright\arraybackslash}p{66mm}@{}}
\toprule
\textbf{Step} & \textbf{Script (directory)} & \textbf{Checkpoint / output} \\
\midrule
Baseline        & \path{few_shot_all.py}, \path{validation_all.py} (\path{all_mixed/}) & \path{ckpt_dir_0_10000/15-6.41}, \path{results_baseline_ood.json} \\
Teacher         & \path{train_teacher.py}, \path{test_teacher.py} (\path{all_mixed/}) & \path{ckpt_dir_0_9972/ckpt-26}, \path{results_teacher.json} \\
Student-scratch & \path{train_scratch.py} (\path{all_mixed/student/}) & \path{ckpt_dirs_scratch/ckpt-08} \\
Student-KD      & \path{train_kd.py} (\path{all_mixed/student/}) & \path{ckpt_dirs_kd_alpha0.5/ckpt-19} \\
Student tests   & \path{evaluate.py} (\path{all_mixed/student/}) & \path{results_student.json} \\
Latency         & \path{inference_benchmark.py} (\path{all_mixed/evaluation/}) & \path{inference_benchmark_results.json} \\
Emulation       & \path{run.py} (\path{mininet_experiment/}) & \path{outputs/six_validation_summary.json} \\
\bottomrule
\end{tabularx}
\end{table}

The emulation is run from the repository root. For example, a model-driven run and a forced-path measurement are started with:
{\footnotesize
\begin{verbatim}
sudo python mininet_experiment/run.py --method kd \
     --seconds 60 --name kd_60_01
sudo python mininet_experiment/run.py --method shortest \
     --candidate 1 --seconds 60 --name path1_new
\end{verbatim}
}
\noindent A model's predictions can be checked without Mininet by running \path{model_worker.py} with the model name and the \path{--check} option.

\paragraph{Notes on the files.}
\begin{itemize}
  \item As committed, \path{validation_all.py} evaluates the original test split and writes \path{results_all.json} (8.22\%). The held-out result in \path{results_baseline_ood.json} (8.40\%) was produced by running the same script with \path{TEST_PATH} set to \path{data/all_mixed_v1/test} and the output file renamed; the console output of that run is in \path{all_mixed/baseline_ood.log}.
  \item The emulation's raw outputs (\path{mininet_experiment/outputs/}, including \path{six_validation_summary.json}, \path{six_validation_predictions.json} and \path{author_real_traffic/}) are excluded from the Git repository. They are kept with the research data backups, and their contents are summarised in \path{mininet_experiment/README.md}.
  \item The console log of the scratch student's training run was not kept; only its checkpoints and evaluation log remain.
\end{itemize}
